\documentclass[11pt]{article}
\usepackage{amsmath,amsthm,amssymb}
\usepackage[margin=1.1in]{geometry}
\usepackage[hidelinks]{hyperref}

\newtheorem{theorem}{Theorem}
\newtheorem{lemma}[theorem]{Lemma}
\theoremstyle{remark}
\newtheorem{remark}[theorem]{Remark}

\title{A Sharp Norm Bound for Weighted Up-Operators on the Boolean Lattice}
\author{An independent contributor}
\date{\today}

\begin{document}
\maketitle

\begin{abstract}
For $a=(a_1,\dots,a_m)\in\mathbb{R}^m$ and $1\le k\le m$, let
$U_k(a):\ell^2\binom{[m]}{k}\to\ell^2\binom{[m]}{k-1}$ be the weighted
up-operator (the weighted $(k,k-1)$ inclusion matrix),
\[
(U_k(a)f)(S)=\sum_{i\notin S} a_i\, f(S\cup\{i\}),\qquad |S|=k-1 .
\]
Writing $\|a\|_2^2=\sum_{i=1}^m a_i^2$, the sharp bound is
\[
\|U_k(a)\|^2\;\le\;\frac{k(m-k+1)}{m}\,\|a\|_2^2
\;=\;\frac{\|U_k(\mathbf 1)\|^2}{m}\,\|a\|_2^2 .
\]
In words: among weight vectors of a fixed Euclidean norm, the
\emph{equal-magnitude} vectors $a=(\pm c,\dots,\pm c)$ maximize the
operator norm, and for interior $k$ and $a\ne0$ they are the only
maximizers. The equal-weight identity $\|U_k(\mathbf 1)\|^2=k(m-k+1)$ is
classical and follows from the Johnson-scheme block diagonalization of
the Boolean lattice; generic weights break that symmetry, so a different
argument is needed. The proof below is a coordinate-deletion induction
on $m$ whose scalar closure defect factors as a perfect square. At the
endpoints $k=1$ and $k=m$ equality holds for every $a$.
\end{abstract}

\section{Introduction}

For integers $m\ge 2$ and $0\le j\le m$ let $\binom{[m]}{j}$ denote the
family of $j$-element subsets of $[m]=\{1,\dots,m\}$, and let
$\ell^2\binom{[m]}{j}$ be the Euclidean space with orthonormal basis
$\{e_S:|S|=j\}$. For a weight vector $a\in\mathbb{R}^m$ and
$1\le k\le m$, the \emph{weighted up-operator}
\[
U_k(a):\ell^2\binom{[m]}{k}\to\ell^2\binom{[m]}{k-1},
\qquad
(U_k(a)f)(S)=\sum_{i\notin S} a_i\, f(S\cup\{i\})
\]
is the $(k,k-1)$ inclusion matrix of the Boolean lattice, with the $i$-th
coordinate scaled by the weight $a_i$. When $a=\mathbf 1=(1,\dots,1)$ this
is the ordinary up-operator of the Boolean lattice; the weight vector
\emph{deforms} that operator while leaving the underlying lattice fixed.

The main result of this note is the sharp inequality
\begin{equation}\label{eq:main}
\boxed{\;\|U_k(a)\|^2\;\le\;\frac{\|U_k(\mathbf 1)\|^2}{m}\,\|a\|_2^2
\;=\;\frac{k(m-k+1)}{m}\sum_{i=1}^m a_i^2\;}
\end{equation}
(Theorem \ref{thm:main}), and the exact characterization of its equality
case. Its content, stated in words, is that \emph{the unweighted operator
is the worst case}: no way of distributing a fixed Euclidean budget of
weight across the $m$ coordinates can make the operator norm exceed what
the flat vector achieves. Equivalently, for fixed $\|a\|_2$ the operator
norm is maximized exactly by the equal-magnitude weights,
$|a_1|=\dots=|a_m|$, and — for interior levels and $a\ne0$ — by nothing
else.

The inequality is sharp because $U_k(\mathbf 1)$ attains the bound. What
makes the statement non-obvious is that it is \emph{not} a corollary of
the classical spectral theory of the Boolean lattice. For $a=\mathbf 1$
the operators $U_k$ and $U_k^{\mathsf T}$ intertwine with the levelwise
action of the symmetric group, so $U_kU_k^{\mathsf T}$ lies in the
Bose--Mesner (Johnson) association scheme and its spectrum is obtained by
a standard simultaneous block diagonalization
\cite{bannai1984,srinivasan2011}. A generic weight vector destroys exactly
that symmetry: the weighted operators no longer lie in the adjacency
algebra of the Johnson scheme, so the equal-weight decomposition says
nothing about $\|U_k(a)\|$ whenever $a$ is not flat. The inequality
nevertheless survives with the \emph{same} optimal constant $C(m,k)=
k(m-k+1)/m$, and the purpose of the proof is to explain why.

The mechanism is a coordinate-deletion induction. Deleting a coordinate
splits $U_k(a)$ into an unweighted block acting on $[m-1]$, a single
weighted block, and the diagonal weight $a_m$ (\eqref{eq:block}); the
operator norm of the resulting $2\times 2$ block operator is then
controlled by a block-norm lemma, and induction on $m$ carries the
constant $C(m,k)$ through unchanged. The induction closes through an
exact scalar inequality whose defect is a perfect square
(\eqref{eq:defect}). That square is the only place where the specific
constant $k(m-k+1)/m$ is pinned down: it is not imposed, it falls out.

We record what the theorem does and does not say. It concerns a
\emph{single} weighted level operator in $\ell^2$-operator norm; it makes
no statement about the full section operator $\sum_i a_i\sigma_i^-$, whose
norm is $\max_k\|U_k(a)\|$, nor about $\ell^p$ norms for $p\ne2$. We make
no claim of an application, and we make no claim that the deletion
induction is the only possible route to \eqref{eq:main}.

\section{Statement of the result}

For $a\in\mathbb{R}^m$ the operator $U_k(a)$ is, in the subset basis, the
weighted $(k,k-1)$ inclusion matrix: its $(S,T)$ entry is $a_i$ if
$T=S\cup\{i\}$ and $0$ otherwise. For $a=\mathbf 1$ its spectral theory is
classical; in particular $\|U_k(\mathbf 1)\|^2=k(m-k+1)$
\cite{bannai1984,srinivasan2011}.

\begin{theorem}\label{thm:main}
For all $m\ge2$, $1\le k\le m$, and $a\in\mathbb{R}^m$,
\[
\|U_k(a)\|^2\;\le\;\frac{k(m-k+1)}{m}\,\sum_{i=1}^m a_i^2 .
\]
The constant $C(m,k)=\frac{k(m-k+1)}{m}$ is best possible. Moreover:
\begin{enumerate}
\item[(i)] for $k=1$ and $k=m$ equality holds for every $a$;
\item[(ii)] for $1<k<m$ and $a\ne0$, equality holds if and only if
$|a_1|=\cdots=|a_m|$; in particular no weight vanishes.
\end{enumerate}
\end{theorem}

The equal-weight case gives
$\|U_k(\mathbf 1)\|^2=k(m-k+1)=C(m,k)\cdot m$, so
\eqref{eq:main} asserts that arbitrary coordinate weights, measured in
$\ell^2$, can never exceed the rescaled classical constant.

\section{Proof}

\subsection{Sign reduction and the two Gram matrices}
For $\sigma\in\{\pm1\}^m$ and $0\le j\le m$ let $D_{\sigma,j}$ be the
diagonal orthogonal operator on $\ell^2\binom{[m]}{j}$ given by
$(D_{\sigma,j}e)_T=\bigl(\prod_{i\in T}\sigma_i\bigr)e_T$. Then, for
$1\le k\le m$,
\begin{equation}\label{eq:sign}
U_k(\sigma\!\cdot\! a)\;=\;D_{\sigma,k-1}\,U_k(a)\,D_{\sigma,k}.
\end{equation}
Indeed, for $T=S\cup\{i\}$ the $(S,T)$ entry of the right side is
$\bigl(\prod_{j\in S}\sigma_j\bigr)a_i\bigl(\prod_{l\in T}\sigma_l\bigr)
=\sigma_i a_i$, the $(S,T)$ entry of the left side; all other entries
agree as zero. Both factors in \eqref{eq:sign} are orthogonal, so
$\|U_k(a)\|=\|U_k(|a|)\|$. We assume $a\ge0$ from now on; the general case
follows, and the equality classification inherits the reduction.

We shall use two Gram matrices of $U_k(a)$. Since $U_k(a)$ is a linear map
between finite-dimensional Euclidean spaces, the \emph{two} Gram
matrices
\[
G_k(a):=U_k(a)^{\mathsf T}U_k(a),
\qquad
H_k(a):=U_k(a)U_k(a)^{\mathsf T}
\]
have the same nonzero eigenvalues --- both are the squares of the
singular values of $U_k(a)$ --- and hence the same largest eigenvalue,
equal by definition to $\|U_k(a)\|^2$. We use $G_k(a)$ where the
levelwise conjugation \eqref{eq:sign} is the natural object (it gives
$G_k(\sigma\!\cdot\!a)=D_{\sigma,k}G_k(a)D_{\sigma,k}$, so the whole
spectrum of $G_k$ is preserved), and $H_k(a)$ in the deletion induction,
where the block structure \eqref{eq:gram} is the natural object.

\subsection{Block decomposition by deleting a coordinate}
Write $a=(a',a_m)$ with $a'\in\mathbb{R}^{m-1}$, and set
$A=U_k(a')$, $B=U_{k-1}(a')$. Order the rows of $U_k(a)$ by $(k-1)$-sets
not containing $m$ first and those containing $m$ second, and the columns
likewise by $k$-sets. Since the $(S,T)$ entry is nonzero only when
$T=S\cup\{i\}$,
\begin{equation}\label{eq:block}
U_k(a)\;=\;
\begin{pmatrix} A & a_m I\\ 0 & B \end{pmatrix},
\end{equation}
where $I$ is the $\binom{m-1}{k-1}\times\binom{m-1}{k-1}$ identity
indexed by the bijection $S\leftrightarrow S\cup\{m\}$ between the
$(k-1)$-sets of $[m-1]$ and the $(k-1)$-sets of $[m]$ containing $m$;
$A$ has shape $\binom{m-1}{k-1}\times\binom{m-1}{k}$ and $B$ has shape
$\binom{m-1}{k-2}\times\binom{m-1}{k-1}$. Consequently
\begin{equation}\label{eq:gram}
H_k(a)=U_k(a)\,U_k(a)^{\mathsf T}\;=\;
\begin{pmatrix} AA^{\mathsf T}+a_m^2 I & a_m B^{\mathsf T}\\
a_m B & BB^{\mathsf T} \end{pmatrix}.
\end{equation}
The off-diagonal blocks have shapes
$\binom{m-1}{k-1}\times\binom{m-1}{k-2}$ and
$\binom{m-1}{k-2}\times\binom{m-1}{k-1}$, matching the block row and
column sizes. (At the boundary level $k=2$ the lower block has
$\binom{m-1}{0}=1$ row; the endpoint levels $k=1$ and $k=m$ are handled
separately below and never require \eqref{eq:block}.)

\subsection{A block-norm lemma}
\begin{lemma}\label{lem:block}
If
\[
M=\begin{pmatrix} P & Q\\ Q^{\mathsf T} & R\end{pmatrix}
\]
is real symmetric with $P,R$ square symmetric, then
\[
\lambda_{\max}(M)\;\le\;
\lambda_{\max}\begin{pmatrix}\lambda_{\max}(P)&s_1(Q)\\ s_1(Q)&\lambda_{\max}(R)\end{pmatrix},
\]
where $s_1$ is the largest singular value.
\end{lemma}
\begin{proof}
For a unit vector $x=(u,v)$, writing $\alpha=|u|$, $\beta=|v|$,
\[
x^{\mathsf T}Mx \le \lambda_{\max}(P)\,\alpha^2+\lambda_{\max}(R)\,\beta^{2}
+2\,s_1(Q)\,\alpha\beta = y^{\mathsf T}Sy,
\qquad y=(\alpha,\beta),\ |y|=1,
\]
by Cauchy--Schwarz applied to $u^{\mathsf T}Qv$. Maximizing over unit $y$
gives $\lambda_{\max}(S)$.
\end{proof}

\subsection{The induction}
We prove the theorem, simultaneously in $k$, by induction on $m$, carrying
the equality clause of Theorem \ref{thm:main}(ii). The endpoints are
exact: $U_1(a)$ is the row vector $(a_1,\dots,a_m)$ and $U_m(a)$ is a
column vector, so
$\|U_1(a)\|^2=\|U_m(a)\|^2=|a|^2=C(m,1)|a|^2=C(m,m)|a|^2$.

Let $1<k<m$ and $a\ne0$. If $a'=0$, then \eqref{eq:block} gives
\[
U_k(a)=\begin{pmatrix} 0 & a_m I\\ 0 & 0\end{pmatrix},
\]
which is $a_m$ times a partial isometry: its only nonzero block is the
identity carrying the $k$-sets of $[m]$ containing $m$ isometrically onto
the $(k-1)$-sets containing $m$, so $\|U_k(a)\|=a_m$ and
$\|U_k(a)\|^2=a_m^2<C(m,k)\,a_m^2$ because
$k(m-k+1)-m=(k-1)(m-k)>0$; the inequality is strict in this case.

Assume then $s:=|a'|^2>0$ and set $t=a_m^2/s$,
$\alpha^2=\|U_k(a')\|^2$, $\beta^2=\|U_{k-1}(a')\|^2$.
Applying Lemma \ref{lem:block} to \eqref{eq:gram} with
$P=AA^{\mathsf T}+a_m^2 I$, $R=BB^{\mathsf T}$, $Q=a_m B^{\mathsf T}$,
and using $s_1(a_m B^{\mathsf T})=a_m\beta$,
\begin{equation}\label{eq:corner0}
\|U_k(a)\|^2\;\le\;\lambda_{\max}
\begin{pmatrix}\alpha^2+a_m^2& a_m\beta\\ a_m\beta&\beta^2\end{pmatrix}.
\end{equation}
By the induction hypothesis applied on $m-1$,
$\alpha^2\le A_0\,s$ and $\beta^2\le B_0\,s$ with
$A_0=C(m-1,k)=\frac{k(m-k)}{m-1}$ and
$B_0=C(m-1,k-1)=\frac{(k-1)(m-k+1)}{m-1}$.
For a symmetric matrix
$\begin{pmatrix} p & o\\ o & q\end{pmatrix}$ the largest eigenvalue
$\tfrac12\bigl(p+q+\sqrt{(p-q)^2+4o^2}\bigr)$ is nondecreasing in $p$,
in $q$, and in $|o|$ (partial derivatives
$\tfrac12\bigl[1\pm(p-q)/\sqrt{(p-q)^2+4o^2}\bigr]\ge0$, and monotone in
$|o|$); hence, from \eqref{eq:corner0},
\begin{equation}\label{eq:corner}
\|U_k(a)\|^2\;\le\; s\cdot
\lambda_{\max}\begin{pmatrix}A_0+t&\sqrt{tB_0}\\ \sqrt{tB_0}&B_0\end{pmatrix}.
\end{equation}

\subsection{Scalar closure}
\begin{lemma}\label{lem:closure}
For $m\ge3$, $1\le k\le m$, $t\ge0$, with $A_0,B_0$ as above and
$C=C(m,k)=\frac{k(m-k+1)}{m}$,
\[
\lambda_{\max}\begin{pmatrix}A_0+t&\sqrt{tB_0}\\ \sqrt{tB_0}&B_0\end{pmatrix}
\;\le\; C(1+t),
\]
and for interior $k$ equality holds if and only if $t=\frac1{m-1}$.
\end{lemma}
\begin{proof}
The claim is $\sqrt{(A_0-B_0+t)^2+4B_0t}\le X(t)$ with
$X(t)=2C(1+t)-A_0-B_0-t$. Now $X$ is affine with slope
$2C-1=\frac{2k(m-k+1)-m}{m}>0$ and intercept
$X(0)=2C-A_0-B_0=\frac{k(k-1)+(m-k)(m-k+1)}{m(m-1)}>0$; hence $X(t)>0$
for $t\ge0$ and squaring is legitimate. The squared defect factors
exactly:
\begin{equation}\label{eq:defect}
X(t)^2-\bigl[(A_0-B_0+t)^2+4B_0t\bigr]
=\frac{4k(k-1)(m-k)(m-k+1)}{m^2(m-1)^2}\,\bigl((m-1)t-1\bigr)^2\ \ge\ 0,
\end{equation}
a polynomial identity in $m,k,t$. For interior $k$ the prefactor is
positive, so equality holds iff $t=\frac1{m-1}$.
\end{proof}

Combining \eqref{eq:corner} with Lemma \ref{lem:closure},
$\|U_k(a)\|^2\le s\,C(1+t)=C(m,k)(s+a_m^2)=C(m,k)|a|^2$, completing the
induction. Sharpness: $G_k(\mathbf 1)\mathbf 1=k(m-k+1)\mathbf 1$, so
equality is attained at $a=\mathbf 1$ in every level. \hfill$\square$

\section{Equality cases}
Throughout this section $1<k<m$ and $a\ne0$. The ``if'' direction: if
$|a_1|=\cdots=|a_m|=c$, the sign reduction \eqref{eq:sign} gives
$\|U_k(a)\|=\|U_k(c\mathbf 1)\|=c\sqrt{k(m-k+1)}=C(m,k)^{1/2}|a|$.

For the only-if direction, suppose equality holds in
Theorem \ref{thm:main}. The chain of \S3 is a chain of upper bounds
\begin{equation}\label{eq:chain}
\|U_k(a)\|^2\;\le\;\text{(corner at }\alpha^2,\beta^2\text{)}
\;\le\;\text{(corner at }A_0s,B_0s\text{)}\;\le\;C(m,k)(1+t)s,
\end{equation}
whose first link is Lemma \ref{lem:block} applied to \eqref{eq:gram},
whose second replaces $(\alpha^2,\beta^2)$ by $(A_0 s,B_0 s)$, and whose
third is Lemma \ref{lem:closure}. Equality in the last term of
\eqref{eq:chain} forces equality in each link.

\emph{Third link.} By Lemma \ref{lem:closure}, equality in
\eqref{eq:defect} forces $t=\frac1{m-1}$; in particular $a_m\ne0$.

\emph{Second link.} Both corners in this comparison are $2\times2$
matrices $\begin{pmatrix} p & o\\ o & q\end{pmatrix}$ with a
\emph{nonnegative} off-diagonal entry: \eqref{eq:corner0} is the corner at
$(p,q,o)=(\alpha^2+a_m^2,\ \beta^2,\ a_m\beta)$, and $s$ times the matrix
of \eqref{eq:corner} is the corner at $(A_0+t,\ B_0,\ \sqrt{tB_0})$; the
link bounds the first by the second through the three monotonicities
recorded before \eqref{eq:corner}. The third link already forces
$t=\frac1{m-1}>0$, hence $a_m\ne0$; and $\beta\ne0$, for $\beta=0$ means
$U_{k-1}(a')=0$, i.e.\ $a'=0$ (each $a_i'$ appears as an entry of
$U_{k-1}(a')$), which is the strict case handled above. Thus both
off-diagonal entries are strictly positive, and for $o>0$ the partial
derivatives
$\tfrac12\bigl[1\pm(p-q)/\sqrt{(p-q)^2+4o^2}\bigr]$ are strictly positive
because $\sqrt{(p-q)^2+4o^2}>|p-q|$ --- the largest eigenvalue is
strictly increasing in $p$ and in $q$ at fixed $o>0$. Equality in the
second link therefore forces equality in the $p$- and $q$-steps, i.e.\
$\alpha^2=A_0 s$ and $\beta^2=B_0 s$.

\emph{Lifting.} The identities $\alpha^2=C(m-1,k)|a'|^2$ and
$\beta^2=C(m-1,k-1)|a'|^2$ are equalities in the induction hypothesis at
$(m-1,k)$ and $(m-1,k-1)$. If $2\le k\le m-2$ the pair $(m-1,k)$ is
interior and the strengthened induction hypothesis applied through the
first identity gives $|a_1|=\cdots=|a_{m-1}|$; this includes $k=2$, since
for $m\ge4$ the pair $(m-1,2)$ is interior. If $k=m-1$ and $m\ge4$, the
pair $(m-1,k)$ is an endpoint and the first identity carries no
information, but $(m-1,m-2)$ is interior and the same conclusion follows
from the second identity. Then $t=\frac1{m-1}$ gives $a_m^2=s/(m-1)$,
the same common value, so $|a_1|=\cdots=|a_m|$, and no weight vanishes.

\emph{Base case $m=3$, $k=2$.} Here $G_2(a)$ is the $3\times3$ matrix
indexed by $\{1,2\},\{1,3\},\{2,3\}$:
\[
G_2(a)=\begin{pmatrix} a_1^2+a_2^2 & a_2a_3 & a_1a_3\\
a_2a_3 & a_1^2+a_3^2 & a_1a_2\\
a_1a_3 & a_1a_2 & a_2^2+a_3^2 \end{pmatrix},
\qquad s=a_1^2+a_2^2+a_3^2,
\]
and $C(3,2)=\frac43$. The certificate $M=\frac43 s\,I-G_2(a)$ is
positive semidefinite: its $1\times1$ principal minors are
$\frac13(a_1^2+a_2^2+4a_3^2)$, $\frac13(a_1^2+4a_2^2+a_3^2)$,
$\frac13(4a_1^2+a_2^2+a_3^2)$; its $2\times2$ principal minors are
$\frac{s}{9}(a_1^2+4a_2^2+4a_3^2)$, $\frac{s}{9}(4a_1^2+a_2^2+4a_3^2)$,
$\frac{s}{9}(4a_1^2+4a_2^2+a_3^2)$; and
\[
\det M=\frac{4}{27}\bigl[(x_1+x_2+x_3)^3-27x_1x_2x_3\bigr],
\qquad x_i=a_i^2,
\]
which is $\ge0$ by the AM--GM inequality
$(x_1+x_2+x_3)^3\ge27x_1x_2x_3$ for $x_i\ge0$, with equality if and only
if $x_1=x_2=x_3$. (The AM--GM equality condition is the standard one:
equality forces $x_1=x_2=x_3$; if some $x_i=0$ the inequality is strict
unless all vanish, which is excluded.) Thus $\lambda_{\max}(G_2(a))\le
\frac43 s$ with equality iff $a_1^2=a_2^2=a_3^2$, and at equality the
nullspace of $M$ is $\operatorname{span}\{(1,1,1)\}$. This completes the
base case and, with the lifting step above, the classification.
\hfill$\square$

\begin{remark}[relation to spectral methods]\label{rem:spectral}
For $a=\mathbf 1$ the operators $\{U_k,U_k^{\mathsf T}\}$ interact with
the rank-grading of the Boolean lattice in a way that makes the
Johnson-scheme diagonalization compute the full spectrum
\cite{srinivasan2011,bannai1984}. For generic weights $a$ the weighted
operators no longer lie in the adjacency algebra of the Johnson scheme,
so that diagonalization does not apply to $U_k(a)$; the induction above
replaces it with a scalar closure whose defect \eqref{eq:defect} is
exactly the deviation from the equal-weight ray. We do not claim that the
deletion induction is the only route to \eqref{eq:main}.
\end{remark}

\section*{Acknowledgments and disclosure}

This note was developed and prepared with the assistance of an automated,
AI-based research agent (an LLM-driven proof-verification and typesetting
pipeline). The block decomposition \eqref{eq:gram}, the scalar-closure
identity \eqref{eq:defect}, and the base-case minors were verified
symbolically or exactly; the companion repository contains reproducible
scripts and a factual development history. Numerical checks additionally
cover random weights for $m\le8$ and every level $k$. \textbf{This note has
not been independently refereed by a human expert.}

\begin{thebibliography}{9}
\bibitem{bannai1984}
E.~Bannai and T.~Ito, \emph{Algebraic Combinatorics I: Association
Schemes}, Benjamin/Cummings, Menlo Park, CA, 1984.
\bibitem{srinivasan2011}
M.~K.~Srinivasan, \emph{Counting spanning trees of the hypercube and its
$q$-analogs by explicit block diagonalization}, arXiv:1104.1481 (2011);
uses the explicit block diagonalization of the Boolean-lattice
(Johnson-scheme) operators.
\end{thebibliography}

\end{document}
